This section defines the test procedures used to evaluate communication performance, closed-loop timing behaviour, simulation fidelity, safety behaviour, and hardware interaction characteristics of the proposed haptic rendering system.
The section focuses on what was tested and how measurements were collected; quantitative outcomes are reported in Section~\ref{sec:results} and interpreted in Section~\ref{sec:discussion}.

\subsection{Communication Performance Tests}

Communication performance was characterised using the dedicated \texttt{serial\_tests} utility described in Section~\ref{sec:logging_communication_validation}. Its rate/integrity and ping--echo modes export CSV files for offline statistical analysis.

Tests were conducted across target transmission rates from 100~Hz to 2000~Hz to assess the effect of increasing packet frequency on reliability and timing performance.

\subsubsection{Packet Integrity and Rate Tests}

Packet integrity and communication timing were evaluated using the rate-test mode of the communication test utility.

For each target transmission frequency, the firmware was configured to stream fixed-format device state packets at a constant rate. Tests were performed at nominal frequencies of 100~Hz, 250~Hz, 500~Hz, 1000~Hz, 1250~Hz, and 2000~Hz.

For each test condition, the host continuously received and parsed incoming packets over the USB virtual COM interface. The recorded performance metrics included valid packet count, dropped packets inferred from sequence gaps, out-of-order packets, host-observed inter-arrival time, effective receive rate, and inter-arrival jitter.

\subsubsection{Round-Trip Latency Test}

Round-trip latency was evaluated using a ping--echo communication test.

In this test, the host transmitted timestamped ping packets to the embedded controller, which immediately returned an echo packet containing the original sequence number and timestamp.

Tests were performed using transmission intervals of 1~ms, 2~ms, and 10~ms. For each test, the evaluated metrics included mean, median, 95th percentile, and 99th percentile RTT, along with minimum and maximum RTT values and timeout counts.

These results were used to characterise communication latency and timing consistency under different transmission loads.

\subsubsection{End-to-End Host-Side Timing Test}

In addition to the standalone communication tests, end-to-end timing of the host-side haptic pipeline was evaluated using the runtime logging framework described in Section~\ref{sec:logging_communication_validation}. Logs captured during normal closed-loop operation were analysed in MATLAB to compute host-side latency, intermediate pipeline delays, matched state-command consistency, parsed state-packet inter-arrival timing, sequence gaps, parsed packet rate, and MCU-side timing trends.

Because the host and microcontroller clocks were not synchronised, the analysis was restricted to host-side timing differences and trend-based interpretation of MCU timestamps. Absolute one-way latency between host and device could therefore not be inferred directly.

\subsection{Simulation Validation} \label{sec:method_simulation}

Simulation validation tests were designed to assess the command-side behaviour of the haptic rendering model independently of physical actuator dynamics.

The simulation tests focused on planar virtual-wall behaviour. A single-depth hold test assessed whether the rendered force command remained steady during sustained penetration. A multi-depth hold test quantified the force--penetration relationship by holding the end-effector at several penetration depths and fitting the logged rendered-force data. Additional free-space and deep-penetration checks verified that the renderer produced negligible force outside contact and that host-side force saturation was applied under extreme penetration. Physical closed-loop interaction with actuator feedback was evaluated separately through the repeated-contact and constrained-motion hardware tests in Sections~\ref{sec:method_stable_contact_hardware} and~\ref{sec:method_constrained_motion_hardware}.

\subsubsection{Virtual Wall Test}

A virtual wall test validated unilateral contact behaviour by moving the end-effector toward and against a planar boundary. The evaluation examined whether force was near zero outside contact, remained steady during a single-depth hold, increased approximately linearly with penetration depth, and matched the implemented spring--damper coupling model while the proxy remained constrained to the wall.

\subsubsection{Free-Space and Saturation Checks}

Free-space and deep-penetration checks were used to verify the boundary cases of the renderer. The free-space check confirmed that no significant force was generated outside contact. The deep-penetration check verified that large virtual penetrations were limited by the configured host-side force clamp before torque commands were generated.

\subsection{Safety Validation}

Safety validation tests evaluated whether the system remained bounded and behaved predictably under both normal and fault conditions. These tests focused on checking torque saturation, rate limiting, and watchdog-based fail-safe behaviour.
The reported safety metrics were saturation fraction, maximum raw and clamped torque, applied torque cap, measured torque slew-rate percentiles, watchdog activation count, total watchdog-active time, activation duration, and sequence-matched command consistency.

\subsubsection{Torque Saturation Test}

A torque saturation test verified that host-side force limits and embedded torque clamps kept actuator commands within safe bounds during high-force contact.

\subsubsection{Rate Limiting Test}

A rate limiting test evaluated the effect of constraining the rate of change of commanded torque at the embedded controller level.

Testing verified that the embedded slew-rate limiter smoothed rapid command changes into bounded torque ramps and reduced high-frequency oscillation during contact transitions.

\subsubsection{Watchdog Timeout Test}

A watchdog timeout test verified safe system behaviour under communication loss by interrupting the incoming command stream during operation.

The test verified that the watchdog triggered when command updates ceased, removed commanded torque within the timeout window, and returned the device to a passive state.

\subsection{Hardware Validation} \label{sec:method_hardware_validation}

Hardware validation assessed the behaviour of the physical actuation system and overall device interaction using the updated runtime logging framework. Because no calibrated force sensor or phase-current measurement was available, the tests did not measure physical output force directly. Instead, logged position, proxy position, rendered force command, host torque command, and device-side applied torque were used as quantitative indicators of contact behaviour, bounded motion, command consistency, and practical interaction quality.

\subsubsection{Repeated Contact Test} \label{sec:method_stable_contact_hardware}

A repeated-contact test evaluated the behaviour of the full closed-loop system during sustained interaction with a virtual surface.

The end-effector was brought into contact with a virtual object while actuator feedback was enabled. Logged data were used to assess whether penetration remained bounded and whether the commanded force--penetration relationship remained consistent with the implemented virtual coupling model. The reported metrics were mean, percentile, and maximum penetration depth, together with the fitted force--penetration stiffness and coefficient of determination. Qualitative observations of oscillation, chatter, and usability were used only to contextualise the logged results.

\subsubsection{Constrained Motion Hardware Test} \label{sec:method_constrained_motion_hardware}

Hardware constrained-motion behaviour was evaluated during physical operation of the device.
The end-effector was brought into contact with a planar virtual surface and moved along it while maintaining contact. Quantitative metrics were computed over a sustained-contact window, excluding entry and release transients. The evaluation assessed whether the logged force command was predominantly normal to the surface, whether normal displacement remained bounded, and whether tangential motion could be produced without large oscillatory excursions. The reported metrics were mean and maximum penetration, normal-to-tangential force ratio, sustained-contact force magnitude, and force coefficient of variation.
